% module_id: physics.quantum.sp1d
% module_version: 1.1.1
% plugin_dir: plugins/quantum_sp1d
% source_files: plugins/quantum_sp1d/__init__.py, plugins/quantum_sp1d/solver.py
% tests: tests/test_quantum_sp1d.py
%
% 本文件是"模块章模板"的样板：后续 8 章必须复刻其节次骨架
%   x.1 模块契约 / x.2 控制方程与物理假设 / x.3 离散化与数值算法 /
%   x.4 耦合求解（仅耦合型模块有此节）/ x.5 验证与基准 /
%   x.6 误差来源与适用边界 / x.7 接口契约与小结
% 任何偏离都视为模板违规。

\chapter{自洽 Schrödinger--Poisson 求解器}
\label{ch:sp1d}

\begin{quote}
\small
\noindent\textbf{模块标识}\quad\pid{physics.quantum.sp1d}\qquad
\textbf{版本}\quad 1.1.1\qquad
\textbf{类别}\quad physics
\end{quote}

\noindent
本章记述二维电子气\index{1concept@概念!2DEG@二维电子气}电荷状态的\emph{自洽}求解：
联立求解静电 Poisson 方程与量子力学 Schrödinger 方程，得到亚带能级、波函数、
占据数以及二维面密度。该模块是 \pid{physics.quantum.analytic_tw}
（\crefchap{ch:analytic_tw}，解析三角阱估计）的即插即用替代实现，二者声明同一
\texttt{quantum} 能力、产出同一组 artifact 键，因此下游模块无需任何改动
\cite{ambacher2000,vurgaftman2003}。

\section{模块契约}
\label{sec:sp1d-contract}\index{4module@模块!physics.quantum.sp1d}\index{6param@参数!algorithm}\index{1concept@概念!plugin architecture@插件架构}

\begin{modulecard}{physics.quantum.sp1d}{1.1.1}
  \cardrow{类别}{physics（物理模型，流水线阶段）}
  \cardrow{能力}{quantum, quantum.sp, quantum.schrodinger\_poisson}
  \cardrow{依赖 artifact}{polarization, materials}
  \cardrow{产出 artifact}{quantum, ns}
  \cardrow{验证状态}{partially-verified}
  \cardrow{源码}{\fpath{plugins/quantum_sp1d/}（多文件插件包）}
  \cardrow{测试}{\fpath{tests/test_quantum_sp1d.py}}
  \cardsep
  \cardrow{算法参数}{algorithm $\in$ \{\akey{newton}, \akey{anderson}\}，默认 \akey{newton}}
  \cardrow{协作者}{经 Deck 的 \akey{nonlinear} 槽解析 \pid{nonlinear.newton}（\crefchap{ch:newton}）}
\end{modulecard}

本模块声明为 \texttt{partially-verified}：其\emph{离散格式}与\emph{导数}均有闭式解
或收敛阶证据支撑（见 \S\ref{sec:sp1d-verification}），但其\emph{绝对数值}依赖
材料常数与屏蔽因子，二者均未经 L2 基准标定（见 \S\ref{sec:sp1d-limits}）。
这一区分贯穿全章：格式正确性可断言，器件数值不可断言。

\section{控制方程与物理假设}
\label{sec:sp1d-model}\index{1concept@概念!conduction band edge@导带边}

\subsection{几何与坐标约定}
\label{sec:sp1d-geometry}\index{1concept@概念!coordinate convention@坐标约定}

求解域为沿外延方向的一维切片，坐标 $z$ 自金属/势垒界面指向体内为正：

\begin{center}
\small
\begin{tabular}{@{}l l l@{}}
\toprule
\textbf{位置} & \textbf{物理含义} & \textbf{边界条件}\\
\midrule
$z=0$      & 金属/AlGaN 界面（肖特基） & Dirichlet：$\Ec(0)=q(\phib-V_g)$\\
$0<z<d$    & AlGaN 势垒（非故意掺杂）  & $\epsr=\epsr^{b}$，$\meff=\meff_{b}$\\
$z=d$      & 异质界面：极化面电荷 + 带阶台阶 & \emph{落在网格面上}（见 \S\ref{sec:sp1d-face}）\\
$d<z<L$    & GaN 沟道与缓冲             & $\epsr=\epsr^{c}$，$\meff=\meff_{c}$\\
$z=L$      & 背面（绝缘缓冲）           & Neumann：零电场\\
\bottomrule
\end{tabular}
\end{center}

竖直方向被解析为量子化方向，横向（栅长方向）不进模型。这是一维近似，
其代价在 \S\ref{sec:sp1d-limits} 中说明。

\subsection{静电方程}
\label{sec:sp1d-poisson}\index{1concept@概念!Poisson equation@Poisson 方程}\index{3method@方法!finite volume@有限体积}

未知量取\emph{导带底} $\Ec(z)$ 而非静电势 $\varphi$。二者相差一个符号
（$\Ec=-q\varphi+\mathrm{const}$），但以 $\Ec$ 为未知量使量子子问题所需的
势剖面\emph{直接可得}，无需在薛定谔求解前再做一次符号转换。定义算子

\begin{equation}
\mathcal{L}[\Ec] \;\equiv\; -\der{}{z}\!\left(\epsr(z)\,\epsO\,\der{\Ec}{z}\right),
\label{eq:poisson-operator}
\end{equation}

则静电方程为

\begin{equation}
\boxed{\;\mathcal{L}[\Ec](z) \;=\; -q\,n(z) \;+\; \sigeff\,\delta(z-d)\;}
\label{eq:poisson}
\end{equation}

其中 $n(z)$ 为体电子密度，$\sigeff$ 为屏蔽后的有效极化面电荷。
\crefeq{eq:poisson} 的符号结构值得强调：源项 $-qn$ 使 $\Ec$ 在电子富集处
\emph{上凸}，即电荷把势阱底部\emph{抹平} —— 这正是自洽求解中屏蔽效应的体现；
界面处 $+\sigeff\delta$ 使 $\Ec$ \emph{向下尖凹}，该凹陷即二维电子气势阱。
若把这两项的符号取反，势阱将变成势垒，2DEG 完全消失。

\crefeq{eq:poisson} 未包含势垒层内的带阶台阶 $\dEc$。带阶是\emph{材料台阶}
而非\emph{电荷}，若将其并入未知量并对它做差分，会在界面处产生一个虚假的
界面偶极子。因此本模块把 $\dEc$ 单独携带，仅在薛定谔方程需要真实势剖面时加上：

\begin{equation}
\Ec^{\mathrm{full}}(z) \;=\; \Ec(z) \;-\; \dEc\,\Theta(z-d),
\label{eq:offset}
\end{equation}

$\Theta$ 为 Heaviside 阶跃。迭代未知量始终是 $\Ec$，即 \crefeq{eq:poisson} 的
\textbf{静电部分}。

\subsection{量子方程}
\label{sec:sp1d-schrodinger}\index{1concept@概念!Schrödinger equation@Schrödinger 方程}\index{3method@方法!BenDaniel--Duke 界面条件}

包络函数近似下，竖直方向的单电子哈密顿量为 BenDaniel--Duke 形式
\cite{benDaniel1966}：

\begin{equation}
\boxed{\;
-\frac{\hb^{2}}{2}\,\der{}{z}\!\left(\frac{1}{\meff(z)}\,\der{\psi_i}{z}\right)
\;+\; \Ec^{\mathrm{full}}(z)\,\psi_i(z)
\;=\; \Ei\,\psi_i(z)
\;}
\label{eq:schrodinger}
\end{equation}

与常见的常质量形式不同，\crefeq{eq:schrodinger} 保留 $1/\meff$ 在微分算子内部。
这是异质界面的正确处理：它给出的界面条件是 $\psi$ 连续、
$(1/\meff)\,\dd\psi/\dd z$ 连续 \cite{benDaniel1966}，而常质量形式会
错误地要求 $\dd\psi/\dd z$ 连续。在 AlGaN/GaN 界面上 $\meff$ 由 $0.30\,\mO$
变到 $0.20\,\mO$，差异不可忽略。

在有效质量近似下，导带被假定为抛物型，只解析 $\Gamma$ 谷；谷简并度 $g_v$
仅作为态密度的倍数进入，不产生额外的能谷结构。

\subsection{占据与面密度}
\label{sec:sp1d-occupation}\index{2qty@物理量!sheet density@面密度}\index{2qty@物理量!Fermi level@费米能级}\index{1concept@概念!subband@亚带}

平衡态下费米能级取为能量零点，$\Ef=0$。第 $i$ 个亚带的二维占据密度由
Fermi--Dirac 统计给出：

\begin{equation}
n_i \;=\; A_i\,\ln\!\left(1+\ee^{(\Ef-\Ei)/\kT}\right),
\qquad
A_i \;=\; \frac{g_v\,\meff_i\,\kT}{\pi\hb^{2}},
\label{eq:occupation}
\end{equation}

其中 $\meff_i$ 为该亚带的态密度有效质量，其定义见 \crefeq{eq:subband-mass}。
总面密度与体密度为

\begin{equation}
\ns \;=\; \sum_i n_i,
\qquad
n(z) \;=\; \sum_i n_i\,\abs{\psi_i(z)}^{2}.
\label{eq:ns}
\end{equation}

束缚态满足 $\Ei<0$（即位于费米能级以下），\crefeq{eq:occupation} 中的
$\ln(1+\ee^{\cdot})$ 形式在 $\Ei$ 远低于 $\Ef$ 时自动退化为线性（简并）区，
无需分段处理；实现中使用 \texttt{logaddexp} 以保证 $\Ei\gg\Ef$ 时不溢出。

由 \crefeq{eq:poisson}，$\sigeff$ 对应的面密度上界为

\begin{equation}
\nsu \;=\; \frac{\sigeff}{q}.
\label{eq:ns-upper}
\end{equation}

$\nsu$ 是物理上界：超过它意味着势垒层内的电场反向，2DEG 不可能存在。
本模块默认\emph{不}施加该钳制（见 \S\ref{sec:sp1d-coupling}）。

\subsection{边界条件与能量零点}
\label{sec:sp1d-bc}\index{1concept@概念!Schottky contact@肖特基接触}

金属/势垒界面为肖特基接触，导带底被钉在

\begin{equation}
\Ec(0) \;=\; q\,(\phib - V_g),
\label{eq:bc-front}
\end{equation}

数值上可用伏特表示。背面 $z=L$ 取零电场（绝缘缓冲）：

\begin{equation}
\der{\Ec}{z}\bigg|_{z=L} \;=\; 0.
\label{eq:bc-back}
\end{equation}

界面的带阶台阶按 \crefeq{eq:offset} 施加。此处有一个易错点：若把 $\dEc$ 的符号
取反，沟道侧导带会被\emph{抬高}而非压低，势阱随之消失。本模块的实现
以势垒侧为参考，沟道侧低 $\dEc$。

\section{离散化}
\label{sec:sp1d-discretisation}\index{3method@方法!discretisation@离散化}

\subsection{面基界面：异质界面是一个"面"，不是一个节点}
\label{sec:sp1d-face}\index{3method@方法!face-based interface@面基界面}

本模块最关键的离散化决策是：\textbf{让异质界面精确落在网格单元之间的面上，
而不是落在某个节点上}。网格节点位于 $d-h/2$ 与 $d+h/2$ 两处，由此

\begin{itemize}[leftmargin=1.7em,itemsep=0.25em,topsep=0.3em]
  \item 导带带阶台阶 $\dEc$、
  \item 介电常数跃变 $\epsr^{b}\to\epsr^{c}$、
  \item 有效质量跃变 $\meff_{b}\to\meff_{c}$、
  \item 极化面电荷 $\sigeff$
\end{itemize}

\noindent
四件事\emph{同时}落在单元边界上。若改为把界面压在节点上，这四件事都会被
以 $O(h)$ 的分辨率抹平，\emph{并且直接吃掉耦合求解的一阶收敛精度} ——
实测结果是从二阶（$p\approx2.08$）退化到一阶（$p\approx1.0$），
详见 \S\ref{sec:sp1d-verification}。因此网格构造函数在必要时会主动拒绝
不合法的输入：当 $d\le h_{\text{fine}}/2$ 时抛出异常，因为此时无法把界面
放到面上。

该决策的另一个后果是：每个单元都完全位于\emph{单一材料}之内，于是中点求积
是二阶的；若界面压在节点上，界面两侧的单元会跨材料，中点求积降为一阶。

\subsection{非均匀网格上的有限体积 Poisson 离散}
\label{sec:sp1d-fv}\index{3method@方法!control volume@控制体}

设节点 $z_1<\dots<z_N$，面间距 $h_i=z_{i+1}-z_i$，单元控制体宽度取

\begin{equation}
w_i \;=\;
\begin{cases}
\tfrac12 h_1, & i=1,\\[2pt]
\tfrac12 (h_{i-1}+h_i), & 1<i<N,\\[2pt]
\tfrac12 h_{N-1}, & i=N,
\end{cases}
\label{eq:cv-width}
\end{equation}

面上的介电系数取两侧的\emph{长度加权串联平均}

\begin{equation}
c_i \;=\; \frac{\epsface_i\,\epsO}{h_i},
\qquad
\epsface_i \;=\; \frac{2}{\dfrac{1}{\epsri{i}}+\dfrac{1}{\epsri{i+1}}}.
\label{eq:face-coeff}
\end{equation}

当材料边界恰好落在该面上时，\crefeq{eq:face-coeff} 正是两个半单元电容的
精确串联组合；在材料均匀处它退化为普通平均。算子 $\mat{L}$ 的第 $i$ 行为

\begin{equation}
(\mat{L}\Ec)_i \;=\;
-c_{i-1}\Eci{i-1} \;+\; (c_{i-1}+c_i)\Eci{i} \;-\; c_i\Eci{i+1},
\qquad 1<i<N,
\label{eq:L-row}
\end{equation}

首行与末行分别为 Dirichlet 与零电场条件：

\begin{equation}
(\mat{L}\Ec)_1 = \Eci{1},
\qquad
(\mat{L}\Ec)_N = -c_{N-1}\Eci{N-1} + c_{N-1}\Eci{N} .
\label{eq:L-bc}
\end{equation}

对分片线性解，该格式是\emph{精确的}（见 \S\ref{sec:sp1d-verification} 中
双介质平行板基准）。

\subsection{界面面电荷的精确消去}
\label{sec:sp1d-sheet}\index{1concept@概念!sheet charge@面电荷}\index{3method@方法!series permittivity average@介电常数串联平均}

面电荷落在面上，因此它不进入 $\mat{L}$，而是被\emph{消去}进右端常数项。
记界面所在面的两侧为节点 $k$（势垒侧）与 $k+1$（沟道侧），界面电位为 $E_I$，
则两个单侧通量为

\begin{equation}
G^{-} = \epsr^{b}\epsO\,\frac{E_I-\Ec^{(k)}}{h/2},
\qquad
G^{+} = \epsr^{c}\epsO\,\frac{\Ec^{(k+1)}-E_I}{h/2},
\qquad
G^{+}-G^{-} = \sigeff .
\label{eq:sheet-flux}
\end{equation}

由 \crefeq{eq:sheet-flux} 解出 $E_I$ 并回代，可验证线性算子 $\mat{L}$ 的形式
\emph{完全不变}，面电荷则被分配为两个节点的源项：

\begin{equation}
b_k \mathrel{+}= \sigeff\,\frac{\epsr^{b}}{\epsr^{b}+\epsr^{c}},
\qquad
b_{k+1} \mathrel{+}= \sigeff\,\frac{\epsr^{c}}{\epsr^{b}+\epsr^{c}} .
\label{eq:sheet-split}
\end{equation}

两个权重之和恒为 $\sigeff$，这是格式自洽的必要条件。常见的"50/50 对分"
只是 $\epsr^{b}=\epsr^{c}$ 的极限情形，在 AlGaN/GaN 上并不成立
（$\epsr^{b}$ 随 Al 组分在 $\sim8.5$–$9.5$ 之间，而 $\epsr^{c}=9.5$），
故本模块采用 \crefeq{eq:sheet-split} 的精确权重。

\subsection{BenDaniel--Duke 算子的对称化}
\label{sec:sp1d-bdd}\index{3method@方法!symmetrisation@对称化}

将 \crefeq{eq:schrodinger} 在控制体上积分，并用 \crefeq{eq:cv-width} 的宽度离散，
得广义本征问题。为把它化为\emph{标准对称}本征问题，作变量替换

\begin{equation}
v_i \;\equiv\; \sqrt{w_i}\,\psi_i ,
\label{eq:symmetrise}
\end{equation}

并定义面系数

\begin{equation}
t_i \;=\; \frac{\hb^{2}}{2\mO}\cdot\frac{\mO}{\bar m_i}\cdot\frac{1}{h_i},
\qquad
\frac{1}{\bar m_i} \;=\; \frac12\!\left(\frac{1}{\meff_i}+\frac{1}{\meff_{i+1}}\right),
\label{eq:t-coeff}
\end{equation}

则对称三对角矩阵 $\mat{A}$ 的两条带为

\begin{equation}
A_{ii} \;=\; \frac{t_{i-1}+t_i}{w_i} \;+\; \Ecifl{i},
\qquad
A_{i,i+1} \;=\; A_{i+1,i} \;=\; -\frac{t_i}{\sqrt{w_i\,w_{i+1}}} .
\label{eq:A-bands}
\end{equation}

于是 \crefeq{eq:schrodinger} 化为 $\mat{A}\mat{v}=\Ei\mat{v}$，可直接调用对称
三对角本征求解器。对称性不是美学选择，而是\emph{后续精确 Jacobian 的前提}：
波函数响应需要对三对角 Hamilton 的正交 resolvent 做一次消元
（\S\ref{sec:sp1d-resolvent}），只有对称形式才能保证该消元的良态与规范选择的
自由度。此外，离散化只在\emph{一个函数}中写定，薛定谔求解与两条响应路径
共用它，因此\emph{导数不可能漂移到它所微分的算子之外}。

界面在有效质量上的处理已由 \crefeq{eq:t-coeff} 的调和平均隐含：
$\bar m_i$ 取两侧的调和平均，等价于 $(1/\meff)\dd\psi/\dd z$ 连续这一
BenDaniel--Duke 界面条件 \cite{benDaniel1966}。

\subsection{子带态密度质量}
\label{sec:sp1d-mass}\index{2qty@物理量!density-of-states mass@态密度有效质量}

\crefeq{eq:occupation} 中的 $\meff_i$ 不是材料的 $\meff(z)$，而是该亚带自身的
态密度有效质量，定义为材料质量在波函数概率密度下的调和平均：

\begin{equation}
\frac{1}{\meff_i} \;=\; \sum_{z} \frac{\abs{\psi_i(z)}^{2}\,w(z)}{\meff(z)}
\qquad\Longleftrightarrow\qquad
\meff_i \;=\; \frac{1}{\sum_z \abs{\psi_i(z)}^{2}w(z)/\meff(z)} .
\label{eq:subband-mass}
\end{equation}

对势垒穿透显著的亚带，\crefeq{eq:subband-mass} 给出的质量明显高于沟道材料值，
这是波函数部分位于势垒层（$\meff_{b}>\meff_{c}$）的直接结果。

\crefeq{eq:subband-mass} 在形式上很不起眼，但它是精确 Jacobian 的一个\emph{必需}
组成项：$\meff_i$ 依赖 $\psi_i$，而 $\psi_i$ 依赖 $\Ec$，故
$\partial n_i/\partial\Ec$ 中含有一条通过 $\meff_i$ 的链。遗漏该项不是
"小近似"，而是一个\emph{缺项}，其误差特征恰好是中心差分最能暴露的那种 ——
\emph{不随差分步长下降的误差地板}，实测约 $1.3\times10^{-5}$
（见 \S\ref{sec:sp1d-jacobian}）。

\section{耦合求解与精确 Jacobian}
\label{sec:sp1d-coupling}\index{3method@方法!self-consistent iteration@自洽迭代}

\subsection{耦合残差与 Newton 系统}
\label{sec:sp1d-residual}\index{3method@方法!Newton--Raphson}\index{6param@参数!tolerance\_eV}

联立 \crefeq{eq:poisson} 与 \crefeq{eq:schrodinger}--\crefeq{eq:ns} 得到一个
关于 $\Ec$ 的非线性代数系统。定义残差

\begin{equation}
\mat{R}(\Ec) \;=\; \mat{L}\Ec \;-\; q\,\mat{W}\,n[\Ec] \;+\; \mat{b},
\label{eq:residual}
\end{equation}

其中 $\mat{W}=\diag(w_i)$（Dirichlet 行置零，因为该行不携带空间电荷），
$\mat{b}$ 汇集了界面面电荷 \crefeq{eq:sheet-split} 与栅极值。
自洽解即 $\mat{R}(\Ec^{\star})=\bm{0}$。

求解策略有二，由参数 \akey{algorithm} 选择：

\begin{itemize}[leftmargin=1.7em,itemsep=0.3em,topsep=0.35em]
  \item \textbf{Newton 法（默认）}：以 $\mat{J}=\partial\mat{R}/\partial\Ec$
        为 Jacobian，交由此 Deck 的 \akey{nonlinear} 槽解析出的求解服务
        （\crefchap{ch:newton}）执行。收敛为二次，实测局部阶 $1.99$。
  \item \textbf{Anderson 加速不动点（回退）}：把 $\Ec\mapsto\mat{L}^{-1}(q\mat{W}n[\Ec]-\mat{b})$
        视作不动点映射，加 Anderson (Pulay) 加速 \cite{anderson1965}。
        不需要协作者模块，用作交叉验证。
\end{itemize}

为何必须用 Newton：裸不动点映射的线性增益约为 $-30$ —— 面密度变化 $1\%$
会使势阱深到让面密度再变化约 $30\%$。纯阻尼迭代要求阻尼因子小于 $0.06$，
并需要数百次迭代；Newton 在个位数迭代内到达同一解。

\subsection{响应算子的三项分解}
\label{sec:sp1d-jacobian}\index{3method@方法!exact Jacobian@精确 Jacobian}\index{3method@方法!perturbation theory@一阶微扰论}

由 \crefeq{eq:residual}，

\begin{equation}
\mat{J} \;=\; \mat{L} \;-\; q\,\mat{W}\,\der{n}{\Ec},
\label{eq:jacobian}
\end{equation}

全部困难集中在 $\dd n/\dd\Ec$。它是\emph{稠密}矩阵（每个节点的势扰动都会
影响每个节点的密度），且由三项构成，三项\emph{全部精确}：

\paragraph{（i）亚带能量响应（主项）。}
势扰动使每个亚带整体平移 $\avg{\psi_i|\delta\Ec|\psi_i}$，从而改变
Fermi--Dirac 占据。由 \crefeq{eq:occupation}，

\begin{equation}
\der{n_i}{\Ei} \;=\; -\frac{A_i}{\kT}\,f_i \;=\; -\beta_i,
\qquad
f_i \;=\; \frac{1}{1+\ee^{\Ei/\kT}},
\label{eq:beta}
\end{equation}

于是（$\avg{\psi_i|\delta\Ec|\psi_i}=\sum_z\abs{\psi_i(z)}^{2}w(z)\,\delta\Ec(z)$）

\begin{equation}
\boxed{\;
\der{n}{\Ec}\bigg|^{(E)}_{z,z'}
\;=\; -\sum_i \beta_i\,\abs{\psi_i(z)}^{2}\,\abs{\psi_i(z')}^{2}\,w(z')
\;}
\label{eq:resp-energy}
\end{equation}

这是一个秩 $n$ 更新（$n$ 为携带的亚带数），承载了主导物理：势阱深度与
电荷之间的负反馈。

\paragraph{（ii）波函数响应。}
势扰动不仅平移能级，还\emph{混合}本征态。一阶微扰论给出

\begin{equation}
\delta\psi_i \;=\; \sum_{j\neq i} \psi_j\,
\frac{\avg{\psi_j|\delta\Ec|\psi_i}}{\Ei-\Ej},
\label{eq:psi-response}
\end{equation}

其对密度的贡献为 $\delta n_i^{(\psi)} = 2 n_i\,\psi_i\,\delta\psi_i$。
\emph{关键之处在于 \crefeq{eq:psi-response} 的求和必须遍历整个谱，不能截断}：
被携带的若干亚带之外，盒态连续谱同样对 $\delta\psi_i$ 有贡献。
旧实现曾把该求和截断到携带的亚带集合内，结果 Jacobian 只能收敛到
局部阶约 $1.5$，残差停在约 $5\times10^{-13}$ 而无法到达机器精度。

\paragraph{（iii）态密度质量响应。}
由 \crefeq{eq:subband-mass}，$\meff_i$ 亦依赖 $\psi_i$，因此

\begin{equation}
\delta\meff_i \;=\; -\,(\meff_i)^{2}\sum_z \frac{2\,\psi_i\,\delta\psi_i\,w(z)}{\meff(z)},
\qquad
\delta n_i^{(m)} \;=\; \frac{n_i}{\meff_i}\,\delta\meff_i .
\label{eq:mass-response}
\end{equation}

该项数量级远小于前两项，却\emph{不可省略}：漏掉它会留下一个不随差分步长
下降的误差地板（实测 $1.3\times10^{-5}$ 相对误差），并使得 Newton 的
末位残差无法压缩到机器精度。

\subsection{正交 resolvent 与秩-1 亏格}
\label{sec:sp1d-resolvent}\index{3method@方法!rank-1 deflation@秩-1 亏格}\index{3method@方法!Sherman--Morrison 公式}\index{3method@方法!pseudo-inverse@伪逆}

\crefeq{eq:psi-response} 的全谱求和在数值上不可直接实施（$O(N^{3})$ 且需要
完整本征分解）。把它写进对称化变量 $v=\sqrt{w}\,\psi$ 后，该求和恰好是
三对角 Hamilton 在携带态之外的正交 resolvent：

\begin{equation}
\sum_{j\neq i}\frac{v_j v_j\T}{\Ei-\Ej}
\;=\; -\sum_{k\neq i}\frac{v_k v_k\T}{E_k-\Ei}
\;=\; -\!\left(\mat{A}-\Ei\mat{I}\right)^{\dagger},
\label{eq:resolvent}
\end{equation}

$\dg$ 记 Moore--Penrose 伪逆。求伪逆的障碍是 $\mat{A}-\Ei\mat{I}$ 奇异
（$v_i$ 张成其零空间）。秩-1 的 Sherman--Morrison 公式
\cite{shermanMorrison1950} 给出

\begin{equation}
\mat{B} = \mat{A}-\Ei\mat{I}+\sigma\,v_i v_i\T
\qquad\Longrightarrow\qquad
\mat{B}^{-1} = \frac{1}{\sigma}v_i v_i\T - \left(\mat{A}-\Ei\mat{I}\right)^{\dagger},
\label{eq:woodbury}
\end{equation}

即伪逆与一个可逆矩阵的逆只差一个秩-1 项。\footnote{秩-$k$ 的推广通常归于
Woodbury 1950 年的普林斯顿技术报告（Memorandum Report 42）；本模块只用到
秩-1 形式，故仅引用 \cite{shermanMorrison1950}。}

本模块\emph{不}选取正则化参数 $\sigma$，而是改用\textbf{固定规范}，从而
完全避开 $\sigma$ 的选择问题：在 $\abs{v_i}$ 最大的节点处摘掉矩阵的行与列
（束缚态零空间向量的每个分量都非零，故任一节点皆可），并把该分量钉为零。
缩减后的三对角块非奇异，其解 $x_g$ 精确满足

\begin{equation}
\left(\mat{A}-\Ei\mat{I}\right)x_g = \bm{b} - (v_i\T\bm{b})\,v_i ,
\label{eq:gauge-system}
\end{equation}

再减去零空间分量即得伪逆作用：

\begin{equation}
\boxed{\;
\left(\mat{A}-\Ei\mat{I}\right)^{\dagger}\bm{b}
\;=\; x_g - (v_i\T x_g)\,v_i
\;}
\label{eq:pseudoinverse}
\end{equation}

\subsection{块右端 Thomas 与复杂度}
\label{sec:sp1d-thomas}\index{3method@方法!Thomas algorithm@Thomas 算法}\index{3method@方法!block right-hand side@块右端}

\crefeq{eq:pseudoinverse} 要求伪逆的\emph{全部 $N$ 列}。关键在于：三对角消元
（Thomas 算法）的消元系数只依赖矩阵本身，与右端无关，因此一次分解即可服务
全部 $N$ 列。本模块把 Thomas 求解器扩展为接受\emph{块右端} $(N\times N)$，
一次前消、$N$ 列同时回代，于是

\begin{equation}
\text{代价} \;=\; O(N^{2})\ \text{每亚带},
\qquad\text{对比全谱本征矢求和的}\ O(N^{3}).
\label{eq:complexity}
\end{equation}

此外，伪逆的右端具有特殊结构：第 $j$ 列为 $\bm{e}_j - v_{i,j}v_i$（规范行置零），
因此无需显式构造单位阵的每一列。

为保证正确性，本模块同时保留一条\emph{独立}实现路径
（\texttt{quantum\_response\_jacobian\_reference}），用 $O(N^{3})$ 的全谱
本征矢求和计算同一个 $\dd n/\dd\Ec$。两条路径共享其余全部装配代码，
\emph{只}在 resolvent 的求法上不同，因此二者的任何分歧都孤立地指向秩-1 亏格
这一处。测试断言二者在舍入误差内一致。

另有一道防御：亏格的前提是 $v_i$ 确实零化 $\mat{A}-\Ei\mat{I}$。若调用者
用手工拼装的过期状态（势剖面与算符不匹配）传入，会得到一个\emph{看起来合理
但毫无意义}的 Jacobian。因此装配前有一道本征对一致性守卫，直接抛出异常
而不是静默继续。

\subsection{热启动与初值}
\label{sec:sp1d-warmstart}\index{3method@方法!warm start@热启动}\index{3method@方法!Fang--Howard 变分形状}\index{3method@方法!bisection@二分法}\index{6param@参数!gamma0\_eV\_m43}

初值只影响收敛速度而不影响收敛结果（有测试证明，见 \S\ref{sec:sp1d-verification}），
但一个接近解的初值对 Newton 的稳健性至关重要。热启动分两步。

\paragraph{第一步：耗尽近似给出面密度估计。}
用二分法求解标量电荷控制关系

\begin{equation}
\begin{aligned}
\ns &= \mathrm{DOS}\cdot\bigl(W(\ns) - E_0(\ns)\bigr),\\[2pt]
W(\ns) &= q\!\left(\sigeff-q\ns\right)\frac{d}{\epsr^{b}\epsO} - q(\phib-\dEc),
\qquad
E_0(\ns)=\gamma_0\,\ns^{2/3},
\end{aligned}
\label{eq:depletion-seed}
\end{equation}

其中三角阱基态项 $E_0=\gamma_0\ns^{2/3}$ 取自 Ambacher 等的电荷控制模型
\cite{ambacher2000}。右端关于 $\ns$ 单调下降，故二分法总能夹住根。
需强调：\crefeq{eq:depletion-seed} 给出的是\emph{解析模型}的答案，
它与自洽解\emph{不同}（同器件下面密度相差 $+4.95\%$，见
\S\ref{sec:sp1d-verification}），此处仅用作起点。

\paragraph{第二步：Fang--Howard 变分形状给出一致剖面。}
以 Fang--Howard 变分波函数形状 \cite{fangHoward1966}

\begin{equation}
\abs{\psi_{\mathrm{FH}}(z)}^{2} \;\propto\; \tfrac12 b^{3}(z-d)^{2}\ee^{-b(z-d)},
\qquad
b \;=\; \left(\frac{33\,\meff_{c}q^{2}\ns}{8\,\epsr^{c}\epsO\hb^{2}}\right)^{1/3},
\label{eq:fang-howard}
\end{equation}

构造初始体密度，做\emph{一次线性} Poisson 求解得到初始 $\Ec$ 剖面。
由于 \crefeq{eq:poisson} 此时已包含界面面电荷，初值本身就带有一个成形的势阱，
Newton 只需修正其形状。

\subsection{Anderson 回退与密度钳制}
\label{sec:sp1d-anderson}\index{3method@方法!Anderson acceleration@Anderson 加速}\index{3method@方法!fixed-point iteration@不动点迭代}\index{6param@参数!density\_clamp}

Anderson 加速器作用在不动点映射 $g$ 上：令残差 $r_k = g(x_k)-x_k$，
以历史增量 $\Delta\mat{X}$、$\Delta\mat{R}$ 求解最小二乘问题
$\gamma = \argmin\norm{\Delta\mat{R}\,\gamma - r_k}_2$，再做

\begin{equation}
x_{k+1} \;=\; x_k + \beta\,r_k - \left(\Delta\mat{X}+\beta\,\Delta\mat{R}\right)\gamma .
\label{eq:anderson}
\end{equation}

当最小二乘解病态（$\norm{\gamma}_\infty$ 超过阈值或非有限）时，清空历史并
退回纯阻尼步。自适应阻尼在残差连续三次变差时把 $\beta$ 减半，下限为
\akey{mixing_min}。

\textbf{密度钳制默认关闭}，这是一个刻意的决定。把 $\ns$ 钳制在
$[0,\nsu]$ 上看似是"物理保险"，但把密度\emph{正好钉在}上界会制造一个
\emph{伪不动点}：两种迭代都可能卡死其上（实测确认）。因此它只作为
最后手段的守卫保留，且\emph{不进入 Jacobian} —— 钳制是不连续的，
对它在边界处求导没有意义。

\section{验证与基准}
\label{sec:sp1d-verification}\index{3method@方法!verification@验证}\index{1concept@概念!convergence order@收敛阶}

本模块的验证策略是分层的：子求解器对\emph{闭式解}，耦合系统对\emph{收敛阶}，
导数对\emph{中心差分}与\emph{独立实现}。所有基准都在
\fpath{tests/test_quantum_sp1d.py} 中以断言形式固化。

\subsection{子求解器对闭式解}

\begin{veriftable}{Poisson 与 Schrödinger 子求解器的闭式解基准}{tab:sp1d-closed-form}
\vgroup{Poisson 算子（\S\ref{sec:sp1d-fv}--\S\ref{sec:sp1d-sheet}）}
\vrow{双介质平行板（无面电荷）}{$<10^{-10}$\U{V}}{$<10^{-10}$\U{V}（分片线性精确）}
\vrow{双介质平行板 $+$ 界面面电荷}{$<10^{-10}$\U{V}}{$<10^{-10}$\U{V}（精确，验证 \crefeq{eq:sheet-split} 的权重）}
\vrow{$\Ec$ 算子 vs.\ $\varphi$ 形式独立离散化}{$<10^{-10}$\U{eV}}{
  $2.9\times10^{-13}$\U{eV}}
\vgroup{Schrödinger 算子（\S\ref{sec:sp1d-bdd}）}
\vrow{无限深方阱 $E_n=n^{2}\pi^{2}\hb^{2}/(2\meff L^{2})$}{二阶收敛}{二阶收敛（相邻误差比 $3$--$5$）}
\vrow{谐振子 $E_n=\hb\omega(n+\tfrac12)$}{$<10^{-3}$\U{eV}}{$<10^{-3}$\U{eV}}
\vrow{三角阱 Airy 零点 $E_n=\abs{\zeta_n}(\hb^{2}F^{2}/2\meff)^{1/3}$}{二阶收敛}{
  $5.5\times10^{-5}\to3.4\times10^{-6}$（加密一档）}
\end{veriftable}

\subsection{导数：精确性而非近似性}

\begin{veriftable}{耦合 Jacobian 的精确性证据}{tab:sp1d-jacobian}
\vrow{三项完整 Jacobian 对中心差分}{相对误差小}{$\sim10^{-11}$（即差分公式本身的噪声水平）}
\vrow{秩-1 亏格 vs.\ 全谱本征矢求和}{$<10^{-9}$}{$4.0\times10^{-13}$}
\vrow{亏格伪逆的定义检验
     $\left(\mat{A}-\Ei\right)\mat{G}=\mat{I}-v_iv_i\T$ 且 $v_i\T\mat{G}=\bm{0}$}{$<10^{-9}$}{通过（逐亚带）}
\vrow{\emph{截断}波函数响应时的误差地板}{---}{
  $1.3\times10^{-5}$，\textbf{不随差分步长下降}（即"缺项"的特征）}
\vnote{最后一行是反例基准：它记录的是\emph{错误}实现的症状，用于说明为何
       \crefeq{eq:mass-response} 与全谱求和不可省略。}
\end{veriftable}

\subsection{耦合系统的收敛阶}

\begin{veriftable}{耦合系统的收敛性与不变量}{tab:sp1d-convergence}
\vrow{\textbf{竖直间距收敛阶}（Richardson 三档）}{$p>1.7$}{\textbf{$p\approx2.08$}
     （面基界面之前为 $\approx1.0$）}
\vrow{\textbf{耦合 Newton 局部收敛阶}}{$>1.9$}{\textbf{$1.99$}
     （三项完整；截断 Jacobian 时仅 $\approx1.5$）}
\vrow{Newton 末位残差}{$<10^{-13}$\U{eV}}{$1.3\times10^{-16}$\U{eV}（机器精度）}
\vrow{Newton vs.\ Anderson 同一答案}{相对差 $<5\times10^{-3}$}{通过（钳制使命中值略偏内）}
\vrow{Newton 迭代数优势}{$\text{it}_{\mathrm{N}}\times5<\text{it}_{\mathrm{A}}$}{
  Newton $10$ 次 vs.\ Anderson $60$ 次，落到同一位数}
\vrow{热启动无关性（$0.35\sigeff/q\to0.95\sigeff/q$）}{相对差 $<5\times10^{-3}$}{$<0.5\%$}
\vrow{求解域深度无关性（$30\to120$\U{nm}）}{相对差 $<10^{-3}$}{$8\times10^{-6}$}
\vrow{竖直加密不变性}{相对差 $<10^{-2}$}{通过}
\vrow{面密度低于极化上界 $\nsu$}{$\ns\le\nsu$}{占据率 $0.83$}
\vrow{温度趋势（$77$ vs.\ $300$\U{K}）}{$0.9\ns^{300}<\ns^{77}\le\ns^{300}$}{通过}
\vrow{负栅压耗尽沟道}{$\ns(-4\,\mathrm{V})<\ns(0\,\mathrm{V})$}{通过}
\end{veriftable}

\noindent
\textbf{域深度无关性}一项值得单独说明：盒态连续谱是有限求解域的产物，
其\emph{单个}波函数并非物理的束缚态。该基准证明它们作为三维连续态的
\emph{求积}是合法的 —— 域深度扩大四倍，面密度只变 $8\times10^{-6}$。

\subsection{与解析模型的对照}

同一个器件（Al$_{0.25}$Ga$_{0.75}$N/GaN，$20$\U{nm} 势垒，
$L_g=150$\U{nm}，屏蔽因子 $0.8$）下，解析三角阱与自洽解的差别具有
明确的\emph{物理形态}：

\begin{veriftable}{自洽解与解析三角阱的对照（同器件）}{tab:sp1d-vs-analytic}
\vrow{量}{解析三角阱}{自洽 SP（差）}
\vrow{$\ns$ (\U{cm^{-2}})}{$7.58287\times10^{12}$}{$7.95848\times10^{12}$（$+4.95\%$）}
\vrow{$\sigeff/q$ 占据率}{$0.7886$}{$0.8276$（$+4.95\%$）}
\vrow{$\Ef$ 距界面带边 (\U{eV})}{$0.461018$}{$0.308840$（$-33.0\%$）}
\vrow{$E_0$ 距界面带边 (\U{eV})}{$0.380340$}{$0.229336$（$-39.7\%$）}
\vrow{$\Ef-E_0$ (\U{eV})}{$0.0806778$}{$0.0795039$（$-1.46\%$）}
\end{veriftable}

\noindent
两个模型对\emph{势}的看法差别很大（$\Ef$ 与 $E_0$ 都掉了约三分之一），
但对\emph{统计}的看法收敛到同一处：决定 $\ns$ 的亚带间距 $\Ef-E_0$
只差 $1.46\%$。这一形态本身就是结论 —— 自洽解给出的势阱更浅，
而其占据由几乎相同的亚带间距支配。

自洽解还给出解析模型根本不具备的量：多亚带阶梯、波函数向势垒的穿透、
由 \crefeq{eq:subband-mass} 决定的真实态密度质量、束缚态数目、
波函数质心与均方根展宽。

\subsection{代价的如实记录}

精确 Jacobian 不是免费的。装配代价比旧的截断版本高约一个数量级：

\begin{veriftable}{精确 Jacobian 的装配代价}{tab:sp1d-cost}\index{1concept@概念!assembly cost@装配代价}
\vrow{网格节点数}{装配（截断 $\to$ 精确）}{整解耗时比}
\vrow{$119$}{$0.6\to5.1$\U{ms}}{$85\to145$\U{ms}（$1.7\times$）}
\vrow{$208$}{$1.3\to16.5$\U{ms}}{$300\to460$\U{ms}（$1.5\times$）}
\vrow{迭代数与线性求解次数}{不变}{不变}
\vnote{计时口径：同机交错计时取 $7$ 次最小值。买到的是一致性
      （\crefchap{ch:transport_dd} 误差预算的前提），\emph{不是}更少的迭代。
      两条后续提速路线见 \S\ref{sec:sp1d-limits}。}
\end{veriftable}

\section{误差来源与适用边界}
\label{sec:sp1d-limits}\index{1concept@概念!validity limits@适用边界}

\limitsheading

\begin{bounds}
  \item \textbf{材料常数未经独立验证。} 极化常数、弹性常数、带隙与有效质量
        转录自文献 \cite{ambacher2000,vurgaftman2003,bernardini1997,dreyer2016}，
        在材料库中标记为 \texttt{unverified}。在通过 Hall 基准（L2）之前，
        \emph{绝对数值不可信}，可断言的是格式与趋势。

  \item \textbf{屏蔽因子没有公认取值。} 参数 \akey{screening}（默认 $0.8$）\index{6param@参数!screening}
        吸收了表面态屏蔽的全部未知物理，是极化链上最大的不确定度来源，
        必须用 Hall 数据标定。

  \item \textbf{有效质量近似，忽略非抛物性。} 导带被假定为抛物型；
        AlGaN/GaN 的强极化使非抛物性不可忽略，尤其是在高面密度下。
        应变对有效质量的修正亦未计入。

  \item \textbf{只有 $\Gamma$ 谷。} 参数 \akey{valley_degeneracy} 只缩放
        态密度，不产生额外的能谷结构。

  \item \textbf{势垒以上的态是有限域盒本征态。} 它们是三维连续态的合法求积
        （域深度依赖 $<10^{-5}$），但\emph{单个}波函数不是物理的束缚态。
        亚带数 \akey{n_subbands} 过少时\index{6param@参数!n\_subbands}，密度将成为势的粗糙泛函，
        拖慢耦合迭代。

  \item \textbf{一维近似。} 横向（短栅效应）、栅边缘场集中、源漏接入区
        均不在模型内。这些属于 \crefchap{ch:transport_dd} 的二维域。

  \item \textbf{Anderson 回退比 Newton 脆弱。} 更慢，且在部分网格上会因
        密度钳制而卡住。默认算法是 Newton。

  \item \textbf{精确 Jacobian 的装配较贵}（见 \creftab{tab:sp1d-cost}）。
        两条正路已明确：(a) 三对角逆的半可分（semiseparable）Green 函数
        两步扫描可把装配降到 $O(N)$ 次向量化运算；(b) 线性步改为
        \emph{矩阵无关}的 Krylov 方法后，响应只需对向量作用
        （每亚带一次三对角求解 $O(N)$），根本不必组装稠密矩阵。

  \item \textbf{无陷阱、无自热、无量热耦合。} 因此电流坍塌、gate-lag、
        drain-lag 一概不在本模块的适用范围。

  \item \textbf{密度钳制默认关闭。} 其理由是数值的而非物理的：钳制会
        制造伪不动点，见 \S\ref{sec:sp1d-anderson}。
\end{bounds}

\section{接口契约与小结}
\label{sec:sp1d-interface}\index{5artifact@数据通道!quantum}\index{5artifact@数据通道!ns}\index{5artifact@数据通道!polarization}\index{1concept@概念!artifact contract@数据通道契约}

\subsection{artifact 契约}

本模块从上下文读取两个 artifact，写回两个。表中列出的字段名与代码中
\texttt{ctx.put}/\texttt{ctx.require} 的实参逐字符一致。

{\small
\begin{longtable}{@{}p{0.06\textwidth} p{0.22\textwidth} p{0.62\textwidth}@{}}
\toprule
\textbf{方向} & \textbf{键} & \textbf{本模块使用/提供的字段}\\
\midrule
\endfirsthead
\toprule
\textbf{方向} & \textbf{键} & \textbf{本模块使用/提供的字段}\\
\midrule
\endhead
\bottomrule
\endfoot
入 & \akey{polarization} &
  \akey{barrier_thickness_nm}, \akey{eps_r_barrier}, \akey{m_eff_barrier},
  \akey{sigma_eff_C_m2}, \akey{x_al}, \akey{band_offset_eV},
  \akey{channel_material}\\
入 & \akey{materials} &
  沟道材料的 \akey{eps_r}, \akey{m_eff}\\
出 & \akey{quantum} &
  契约键 \akey{ns_m2}, \akey{ns_cm2}, \akey{ef_eV}, \akey{e0_eV},
  \akey{phi_b_eV}, \akey{delta_ec_eV}, \akey{converged}, \akey{iterations}；
  另加亚带能级/质量/分密度、束缚态数、质心与展宽、残差历史、剖面等自洽细节\\
出 & \akey{ns} & 标量面密度（\U{cm^{-2}}），供提取器直接消费\\
\end{longtable}
}

契约键中的前七项与 \pid{physics.quantum.analytic_tw} 完全一致，这是
\emph{即插即用}的机械保证：整个替换只需改动 Deck 中的一行模块绑定，
\pid{extract.fom}（\crefchap{ch:fom}）与流水线均无须改动，因为下游只认
artifact 键而不认实现。

\subsection{小结}

本章的核心结论可以概括为三点。

其一，\textbf{异质界面必须作为"面"处理}。把带阶、介电跃变、质量跃变与面电荷
四件事同时放到单元边界上，耦合求解才恢复二阶收敛（$p\approx2.08$）；
压在节点上会同时损失精度与一个收敛阶。

其二，\textbf{Jacobian 的精确性是可达的，但需要把三项都写对}。
亚带能量响应、波函数响应（全谱、经秩-1 亏格降为 $O(N^{2})$）与态密度质量响应
三者缺一不可。付清这笔账之后，Newton 的局部阶回到 $1.99$，
末位残差落到机器精度 $1.3\times10^{-16}$。

其三，\textbf{格式可信与数值可信是两件事}。本模块的离散格式有闭式解与收敛阶
背书，但绝对数值受制于未标定的材料常数与屏蔽因子。这一区分在本文档中
始终显式给出，不作合并陈述。

\cleardoublepage
